\section{Discussion and Conclusions}
\label{sec:discussion}

\subsection{RQ1}

\emph{How do relative speedup and efficiency of the same parallel code
evolve with the number of threads, natively and in WASM in the browser?}

With enough work per parallel region, the same way. On a loop with no
synchronization the efficiency falls slowly with $N$, to 0.87 at eight
threads natively and in Chrome and 0.92 in Firefox; Chrome's speedup is
within about 1\% of native's, and the browser over native time ratio does
not grow with $N$, in strong and in weak scaling. A loop that does not
touch shared memory scales almost ideally in the three environments (the
local histogram, 0.98 to 1.06).

WASM costs a constant factor per thread, 1.17 in Chrome and 1.36 in Firefox
on the filter, measured here as a third more instructions per pixel run at
different rates. That factor divides with the work. What does not divide is
the fixed cost of a parallel region, six to ten times larger in the
browsers: negligible for regions of milliseconds, half of the time for
regions of about 0.4~ms at eight threads.

The one significant RQ1 difference, Firefox's higher speedup, is not
evidence that a browser runs threads better. The sub-linear speedup of the
filter is not a serial fraction (the timed region has none) but a cost of
eight threads sharing memory, and a slower loop asks less of the memory:
across three builds and three environments, the faster one thread reads
memory the less eight threads gain. Firefox is the slowest environment at
every thread count. This explanation is supported by counters, by the SIMD
builds, natively by the image size, and by one intervention inside each
engine; none of them intervenes on the memory itself, and about half of the
native loss is time outside the loop that was not measured.

\subsection{RQ2}

\emph{To what extent does synchronization between threads affect the
scalability gap between WASM and native?}

It depends on the engine, and on how the contention is spread. In Chrome,
relative to the lock-free control, shared atomic counters cost 1.4 to 2.3
times what they cost natively when contention is concentrated, and 0.42
times when it is spread over 256 counters. In Firefox they cost what they
cost natively in absolute time; its ratios below 1 come from a slower
control, not from cheaper synchronization.

The difference is a compiler's choice, not a property of WASM threads.
Memory ordering is not the reason: WASM atomics are sequentially
consistent~\cite{watt2019}, but on x86 an atomic add costs the same
whatever ordering is asked~\cite{schweizer2015}. V8's optimizing compiler
turns an atomic add whose result is unused into a compare-and-swap loop.
That sequence is slower uncontended, is retried once in four adds under
concentrated contention, and is faster when eight cores spread over 256
lines. A native program with the same sequence reproduces Chrome's times in
both directions; the same engine with its baseline compiler, and the other
engine, show no penalty at $k=1$. One source and one set of flags do not
give one instruction stream, so a statement about ``WASM atomics'' is a
statement about an engine version.

\subsection{For developers}

\begin{itemize}[leftmargin=1.5em]
\item Porting a data-parallel loop to WASM threads keeps its scaling and
costs a constant factor; expect 1.2 to 1.4 times the native time on loops
like these.
\item Keep parallel regions coarse. At eight threads a region needs about
0.4 to 0.5~ms of sequential work in a browser to get half the ideal
speedup, three to four times what native code needs, and more if it is
called once and cold.
\item Inside a region (Section~\ref{sec:grain}), Rayon's adaptive split is within 3\% of the best
fixed grain at these sizes in the three environments; there is no reason
to tune the grain per environment.
\item Avoid shared counters, as natively: they do not scale anywhere.
Privatize and merge. Where atomics are unavoidable, give each its own cache
line, and do not assume that their cost in one engine predicts another.
\end{itemize}

\subsection{For benchmarkers}

\begin{itemize}[leftmargin=1.5em]
\item A CPU mask, a fresh profile and a pinned version do not by
themselves confine, quiet and freeze a browser. Record the state with every
run and reject runs by rules fixed in advance.
\item Read the machine code. Both findings of this study, and two of our
own mistakes, were visible only there.
\item A relative speedup needs its baseline: a slower environment can
scale better because it is slower.
\item Whole-run hardware counters of a browser are mostly not the code
under test.
\item Write the prediction before the test. One of ours failed in part,
and it is the most informative line of Section~\ref{sec:why-rq2}.
\end{itemize}

\subsection{What stays open}

The mechanism that makes a compare-and-swap loop faster than a locked add
when eight cores spread over 256 padded counters. Why SpiderMonkey's filter
loop, as long as V8's, runs fewer instructions per cycle, and why its local
histogram misses more caches with one thread. What Chrome loses outside the
loop at eight threads, and whether the idle and wake path of the WASM
thread pool is what makes the browsers' fork-join slower. The efficiency
dip at six and seven threads. What sets the state of contention of a run
at $k=1$. The slow early iterations of one Firefox cell. And everything
outside this design: other machines, turbo, other engine versions, locks
and irregular parallelism.

\subsection{Conclusion}

The same Rust/Rayon code scales in the browser as natively when each
parallel region has enough work: WASM threads cost a roughly constant
factor and a slower fork-join, and where a browser scales better it is
because its code is slower on a memory-bound loop. Shared atomics widen the
gap in Chrome under concentrated contention and narrow it under spread
contention; Firefox is never worse than native. The cause is V8's
compare-and-swap loop for an atomic add, reproduced natively and absent
from V8's other compiler. The study was measured twice in full; the second
measurement, with its corrections and a rule for adopting it fixed in
advance, gives the same ten decisions as the first.
